\documentclass[a4paper]{article}
\usepackage[fontset=fandol]{ctex}
\usepackage{amsmath, amssymb}
\usepackage{graphicx}
\usepackage[margin=2.5cm]{geometry}
\usepackage[most]{tcolorbox}
\usepackage{etoolbox}
\usepackage{listings}
\usepackage{hyperref}
\usepackage{booktabs}
\usepackage{subcaption}
\usepackage{float}
\usepackage{tikz}
\IfFileExists{pgfplots.sty}{\usepackage{pgfplots}\pgfplotsset{compat=1.18}}{}
\newtcolorbox{knowledgebox}[1]{enhanced, colback=blue!5!white, colframe=blue!75!black, colbacktitle=blue!75!black, coltitle=white, fonttitle=\bfseries, title=#1, attach boxed title to top left={yshift=-2mm, xshift=2mm}, boxrule=1pt, sharp corners}
\newtcolorbox{importantbox}[1]{enhanced, colback=yellow!10!white, colframe=yellow!80!black, colbacktitle=yellow!80!black, coltitle=black, fonttitle=\bfseries, title=#1, sharp corners}
\newtcolorbox{warningbox}[1]{enhanced, colback=red!5!white, colframe=red!75!black, colbacktitle=red!75!black, coltitle=white, fonttitle=\bfseries, title=#1, sharp corners}
\lstset{language=Python,basicstyle=\ttfamily\small,keywordstyle=\color{blue},stringstyle=\color{red!60!black},commentstyle=\color{green!60!black},breaklines=true,frame=single,numbers=left,numberstyle=\tiny\color{gray},captionpos=b,extendedchars=false}
\newcommand{\notetitle}{Feature Extraction}
\newcommand{\noteauthors}{Codex 基于课程页、公开视频字幕、coverage ledger 与补充官方材料整理}
\newcommand{\notedate}{\today}
\newcommand{\videochannel}{WAVLab}
\newcommand{\videopublishdate}{2023-09-11}
\newcommand{\videoduration}{unknown}
\newcommand{\videourl}{https://www.youtube.com/watch?v=Z5v\_hj\_\_CWU}
\newcommand{\videocoverpath}{cover.jpg}
\begin{document}
\begin{titlepage}
\centering
{\Large 课程笔记\par}
\vspace{1.2cm}
{\huge\bfseries \notetitle\par}
\vspace{0.8cm}
{\large \noteauthors\par}
\vspace{0.3cm}
{\large \notedate\par}
\vspace{1.2cm}
\ifdefempty{\videocoverpath}{{\small 公开视频封面不可用。\par}}{\includegraphics[width=0.82\textwidth,height=0.45\textheight,keepaspectratio]{\videocoverpath}\par}
\vfill
\begin{tcolorbox}[width=0.9\textwidth, colback=black!2!white, colframe=black!60, sharp corners]
\textbf{视频作者/频道}：\videochannel\par
\textbf{发布日期}：\videopublishdate\par
\textbf{视频时长}：\videoduration\par
\textbf{视频链接}：\href{\videourl}{\nolinkurl{\videourl}}
\end{tcolorbox}
\end{titlepage}
\tableofcontents
\newpage
\section{4.0 课程主题总览}
根据课程页主题、公开视频字幕以及本讲的 coverage ledger，4.0 课程主题总览 这一部分主要围绕 Feature Extraction | Basic pipeline | Some advances in feature extractions 展开。
写作上保持 coverage-first 原则：先把老师在这一阶段真正讲过的问题设定、模型部件、训练或推理逻辑讲清，再压缩成适合复习的结构。

\begin{figure}[H]
\centering
\includegraphics[width=0.92\textwidth]{figures/coverage_map.png}
\caption{本讲 coverage map：按 coverage\_units.jsonl 归纳出的核心主题与章节映射。}
\end{figure}

\begin{figure}[H]
\centering
\includegraphics[width=0.92\textwidth]{figures/timeline.png}
\caption{本讲 timeline：根据字幕时间范围归纳的主题推进顺序。}
\end{figure}

\subsection{Feature Extraction}
从证据位置看，这一部分对应的课程证据区间是 2023-09-11。课程在这里强调的核心内容可以概括为：Feature Extraction | Basic pipeline | Some advances in feature extractions。
这部分相当于把波形如何变成可建模特征的整条前端链路重新拆开，方便读者理解后续 acoustic model 到底在消费什么信息。

在本章结构中，这一单元被并入“4.0 课程主题总览”，目的是把同一阶段的定义、机制和权衡集中讲清。

\subsection{本章小结}
4.0 课程主题总览 这一部分把 Feature Extraction 中相邻的教学单元串成了一条连续叙事：既交代了这一阶段要解决的问题，也交代了方法、推理或工程权衡为什么要这样组织。

\section{4.1 语音信号与采样}
根据课程页主题、公开视频字幕以及本讲的 coverage ledger，4.1 语音信号与采样 这一部分主要围绕 Feature extraction pipeline | Signal acquisition | Preprocessing motivation 展开。
写作上保持 coverage-first 原则：先把老师在这一阶段真正讲过的问题设定、模型部件、训练或推理逻辑讲清，再压缩成适合复习的结构。

\subsection{Feature extraction pipeline}
从证据位置看，这一部分对应的课程证据区间是 00:00:00,919 -- 00:06:38,039。课程在这里强调的核心内容可以概括为：Feature extraction pipeline | Signal acquisition | Preprocessing motivation。
这部分相当于把波形如何变成可建模特征的整条前端链路重新拆开，方便读者理解后续 acoustic model 到底在消费什么信息。

在本章结构中，这一单元被并入“4.1 语音信号与采样”，目的是把同一阶段的定义、机制和权衡集中讲清。

\subsection{本章小结}
4.1 语音信号与采样 这一部分把 Feature Extraction 中相邻的教学单元串成了一条连续叙事：既交代了这一阶段要解决的问题，也交代了方法、推理或工程权衡为什么要这样组织。

\section{4.2 预处理、downsampling 与 framing}
根据课程页主题、公开视频字幕以及本讲的 coverage ledger，4.2 预处理、downsampling 与 framing 这一部分主要围绕 Sampling | Downsampling | Preprocessing and framing entry point 展开。
写作上保持 coverage-first 原则：先把老师在这一阶段真正讲过的问题设定、模型部件、训练或推理逻辑讲清，再压缩成适合复习的结构。

\subsection{Sampling}
从证据位置看，这一部分对应的课程证据区间是 00:06:33,680 -- 00:13:37,839。课程在这里强调的核心内容可以概括为：Sampling | Downsampling | Preprocessing and framing entry point。
从课程组织上看，这一部分承担的是把局部概念嵌回整条识别流水线中的作用，避免读者把公式或模块孤立地记忆。

在本章结构中，这一单元被并入“4.2 预处理、downsampling 与 framing”，目的是把同一阶段的定义、机制和权衡集中讲清。

\subsection{本章小结}
4.2 预处理、downsampling 与 framing 这一部分把 Feature Extraction 中相邻的教学单元串成了一条连续叙事：既交代了这一阶段要解决的问题，也交代了方法、推理或工程权衡为什么要这样组织。

\section{4.3 STFT / spectrogram}
根据课程页主题、公开视频字幕以及本讲的 coverage ledger，4.3 STFT / spectrogram 这一部分主要围绕 Framing | Windowing | Short-time spectral analysis 展开。
写作上保持 coverage-first 原则：先把老师在这一阶段真正讲过的问题设定、模型部件、训练或推理逻辑讲清，再压缩成适合复习的结构。

\subsection{Framing}
从证据位置看，这一部分对应的课程证据区间是 00:13:34,519 -- 00:21:06,360。课程在这里强调的核心内容可以概括为：Framing | Windowing | Short-time spectral analysis。
从课程组织上看，这一部分承担的是把局部概念嵌回整条识别流水线中的作用，避免读者把公式或模块孤立地记忆。

在本章结构中，这一单元被并入“4.3 STFT / spectrogram”，目的是把同一阶段的定义、机制和权衡集中讲清。

\subsection{本章小结}
4.3 STFT / spectrogram 这一部分把 Feature Extraction 中相邻的教学单元串成了一条连续叙事：既交代了这一阶段要解决的问题，也交代了方法、推理或工程权衡为什么要这样组织。

\section{4.4 Mel scale / filterbanks}
根据课程页主题、公开视频字幕以及本讲的 coverage ledger，4.4 Mel scale / filterbanks 这一部分主要围绕 Spectrograms | Mel scale | Filterbanks 展开。
写作上保持 coverage-first 原则：先把老师在这一阶段真正讲过的问题设定、模型部件、训练或推理逻辑讲清，再压缩成适合复习的结构。

\subsection{Spectrograms}
从证据位置看，这一部分对应的课程证据区间是 00:21:02,440 -- 00:27:50,279。课程在这里强调的核心内容可以概括为：Spectrograms | Mel scale | Filterbanks。
这部分相当于把波形如何变成可建模特征的整条前端链路重新拆开，方便读者理解后续 acoustic model 到底在消费什么信息。

在本章结构中，这一单元被并入“4.4 Mel scale / filterbanks”，目的是把同一阶段的定义、机制和权衡集中讲清。

\subsection{本章小结}
4.4 Mel scale / filterbanks 这一部分把 Feature Extraction 中相邻的教学单元串成了一条连续叙事：既交代了这一阶段要解决的问题，也交代了方法、推理或工程权衡为什么要这样组织。

\section{4.5 MFCC / cepstrum}
根据课程页主题、公开视频字幕以及本讲的 coverage ledger，4.5 MFCC / cepstrum 这一部分主要围绕 MFCC | Cepstrum | DCT and compact acoustic features 展开。
写作上保持 coverage-first 原则：先把老师在这一阶段真正讲过的问题设定、模型部件、训练或推理逻辑讲清，再压缩成适合复习的结构。

\subsection{MFCC}
从证据位置看，这一部分对应的课程证据区间是 00:27:46,360 -- 00:34:44,480。课程在这里强调的核心内容可以概括为：MFCC | Cepstrum | DCT and compact acoustic features。
这部分相当于把波形如何变成可建模特征的整条前端链路重新拆开，方便读者理解后续 acoustic model 到底在消费什么信息。

在本章结构中，这一单元被并入“4.5 MFCC / cepstrum”，目的是把同一阶段的定义、机制和权衡集中讲清。

\subsection{本章小结}
4.5 MFCC / cepstrum 这一部分把 Feature Extraction 中相邻的教学单元串成了一条连续叙事：既交代了这一阶段要解决的问题，也交代了方法、推理或工程权衡为什么要这样组织。

\section{4.6 normalization / delta / context}
根据课程页主题、公开视频字幕以及本讲的 coverage ledger，4.6 normalization / delta / context 这一部分主要围绕 Normalization | Delta features | Context stacking 展开。
写作上保持 coverage-first 原则：先把老师在这一阶段真正讲过的问题设定、模型部件、训练或推理逻辑讲清，再压缩成适合复习的结构。

\subsection{Normalization}
从证据位置看，这一部分对应的课程证据区间是 00:34:42,040 -- 00:41:28,560。课程在这里强调的核心内容可以概括为：Normalization | Delta features | Context stacking。
这部分相当于把波形如何变成可建模特征的整条前端链路重新拆开，方便读者理解后续 acoustic model 到底在消费什么信息。

在本章结构中，这一单元被并入“4.6 normalization / delta / context”，目的是把同一阶段的定义、机制和权衡集中讲清。

\subsection{本章小结}
4.6 normalization / delta / context 这一部分把 Feature Extraction 中相邻的教学单元串成了一条连续叙事：既交代了这一阶段要解决的问题，也交代了方法、推理或工程权衡为什么要这样组织。

\section{4.7 learned front ends}
根据课程页主题、公开视频字幕以及本讲的 coverage ledger，4.7 learned front ends 这一部分主要围绕 Advanced front ends | Learned features | Robustness 展开。
写作上保持 coverage-first 原则：先把老师在这一阶段真正讲过的问题设定、模型部件、训练或推理逻辑讲清，再压缩成适合复习的结构。

\subsection{Advanced front ends}
从证据位置看，这一部分对应的课程证据区间是 00:41:25,160 -- 00:47:28,640。课程在这里强调的核心内容可以概括为：Advanced front ends | Learned features | Robustness。
这部分相当于把波形如何变成可建模特征的整条前端链路重新拆开，方便读者理解后续 acoustic model 到底在消费什么信息。

在本章结构中，这一单元被并入“4.7 learned front ends”，目的是把同一阶段的定义、机制和权衡集中讲清。

\subsection{本章小结}
4.7 learned front ends 这一部分把 Feature Extraction 中相邻的教学单元串成了一条连续叙事：既交代了这一阶段要解决的问题，也交代了方法、推理或工程权衡为什么要这样组织。

\section{4.8 practical tradeoffs}
根据课程页主题、公开视频字幕以及本讲的 coverage ledger，4.8 practical tradeoffs 这一部分主要围绕 Practical tradeoffs | Why feature design still matters 展开。
写作上保持 coverage-first 原则：先把老师在这一阶段真正讲过的问题设定、模型部件、训练或推理逻辑讲清，再压缩成适合复习的结构。

\subsection{Practical tradeoffs}
从证据位置看，这一部分对应的课程证据区间是 00:47:26,599 -- 00:53:28,000。课程在这里强调的核心内容可以概括为：Practical tradeoffs | Why feature design still matters。
这部分相当于把波形如何变成可建模特征的整条前端链路重新拆开，方便读者理解后续 acoustic model 到底在消费什么信息。

在本章结构中，这一单元被并入“4.8 practical tradeoffs”，目的是把同一阶段的定义、机制和权衡集中讲清。

\subsection{本章小结}
4.8 practical tradeoffs 这一部分把 Feature Extraction 中相邻的教学单元串成了一条连续叙事：既交代了这一阶段要解决的问题，也交代了方法、推理或工程权衡为什么要这样组织。

\section{总结与延伸}
Feature Extraction 这一讲在整门《Speech Recognition and Understanding》课程中的位置，是把 Feature Extraction 讲成一个可复用的建模模块。
从教材化角度看，本讲最重要的不是记住单一术语，而是理解它如何嵌入完整的 automatic speech recognition pipeline：输入语音如何被表示、模型如何被训练、解码如何得到最终转写，以及这些设计分别带来什么假设和权衡。
由于公开源以课程页和公开视频字幕为主、缺少官方 slide PDF，本章在正文里尽量把术语、方法和工程语境解释完整，并把缺失项留在 omission 和 manifest 里，而不是凭空补写。
\end{document}
